O presente relatório de estágio descreve o desenvolvimento de uma solução automatizada para a empresa Veeries, organização que atua no setor do agronegócio oferecendo produtos de análise de mercado e fluxo de \textit{commodities}. O problema central abordado foi o severo gargalo operacional na ingestão de dados oriundos de relatórios em formato PDF de duas fontes externas confidenciais. Tal extração era executada de forma estritamente manual pela equipe de analistas, consumindo cerca de três horas por lote e apresentando alta suscetibilidade a falhas humanas e estruturais. Para solucionar este desafio tecnológico e de negócio, foi projetado e implementado um \textit{pipeline} de dados orientado a eventos, fundamentado na Arquitetura \textit{Medallion} (composto pelas camadas progressivas \textit{Bronze, Silver} e \textit{Gold}) e construído na linguagem Python, com suporte das bibliotecas Pandas e Pandera. O fluxo automatizou a extração bruta do formato não estruturado e implementou rigorosos Contratos de Dados (\textit{schemas}) para validações matemáticas e semânticas, além de adotar uma estratégia arquitetural de interrupção imediata em caso de anomalias estruturais (\textit{fail-fast}). Como resultado direto da implantação, o tempo de processamento de um lote de relatórios foi reduzido de horas para um intervalo de 2 a 5 segundos. A esteira automatizada transferiu a responsabilidade técnica integralmente para a equipe de desenvolvedores, livrou os analistas de mercado de tarefas braçais de transcrição e conferiu alta governança, rastreabilidade e resiliência ao banco de dados da organização. O sistema provou-se capaz de interceptar e bloquear erros de publicação oriundos das próprias fontes emissoras antes que alcançassem as ferramentas analíticas e de tomada de decisão da diretoria.

\vspace{\onelineskip}
\noindent
\textbf{Palavras-chave:} Engenharia de Dados. Arquitetura \textit{Medallion}. \textit{Pipeline} de Dados. Qualidade de Dados. Automação.